% ============================================================
% 5. EXPERIMENTS
% ============================================================
\section{Experiments}
\label{sec:experiments}

\subsection{Experimental Setup}
\label{sec:setup}

We evaluate EAAG along three axes: (1)~Does environment-aware gating
Pareto-dominate fixed-direction baselines across diverse environments?
(\S\ref{sec:main-results}) (2)~Which components (direction learning,
multi-signal features, LLM reasoning) drive the gains?
(\S\ref{sec:ablation}) (3)~Do the Two-Source Model's predictions
hold empirically? (\S\ref{sec:theory-verification})

\paragraph{Environments.} We evaluate on 8 environments spanning QA,
code generation, web navigation, fact verification, text games, and
crafting planning (Table~\ref{tab:env-setup}): HotpotQA,
APPS Intro, WebShop, FEVER, TWExpress, Plancraft, APPS Interview,
and CRUXEval. These environments are selected to cover the full
range of the Two-Source Model, from pure Type~I (FEVER) through
mixed ($\approx p_I^*$, APPS Intro) to Type~D (APPS Interview),
including extreme rollout properties (TWExpress: rollout-safe;
Plancraft: rollout-harmful).

\paragraph{Baselines.} All methods share the same agent and
optimizer $T$; we compare only the gate decision.
(1)~\emph{Bounds}: base\_only, always\_trigger, oracle.
(2)~\emph{Fixed-direction CB}: CaTS, SEAG, CoRefine, CATTS, AUQ,
s1\_budget, all assuming high uncertainty $\to$ trigger.
(3)~\emph{Ablation}: BSW (intentionally reversed direction).
(4)~\emph{Ours}: EAAG.
Full environment specifications and hyperparameter details are
provided in Appendix~\ref{app:experiments}.

\paragraph{Metrics.} SR (success rate, $\uparrow$) and Cost
(rollouts per episode during deployment, $\downarrow$).
A method Pareto-dominates another if $\text{SR} \geq$ and
$\text{Cost} \leq$ with at least one strict inequality.
Cost measures only the \emph{deployment} rollout budget; calibration
or exploration overhead is excluded for all methods, ensuring a fair
comparison of runtime efficiency.

\begin{table}[t]
\caption{Environment setup. Base/Always: SR without/with optimizer.
$\Delta$: rollout headroom.}
\label{tab:env-setup}
\centering\small
\begin{tabular}{lcccl}
\toprule
Environment & Base SR & Always SR & $\Delta$ & Optimizer $T$ \\
\midrule
HotpotQA      & 49.0\% & 97.0\% & +48.0 & Per-action eval \\
APPS Intro    & 58.5\% & 64.5\% &  +6.0 & $K$-variant sampling \\
WebShop       &  7.2\% & 43.0\% & +35.8 & LLM-Propose-$K$ \\
FEVER         & 37.0\% & 99.8\% & +62.8 & Per-action eval \\
TWExpress     & 67.5\% & 99.3\% & +31.8 & Per-action eval \\
Plancraft     & 29.8\% & 22.8\% &  $-$7.0 & Per-action eval \\
APPS Interview& 60.5\% & 79.5\% & +19.0 & $K$-variant sampling \\
CRUXEval      & 85.0\% & 99.5\% & +14.5 & $K$-variant sampling \\
\bottomrule
\end{tabular}
\end{table}


\subsection{Main Results}
\label{sec:main-results}

\textbf{EAAG outperforms fixed-direction baselines across
environments and backbones.} As shown in Table~\ref{tab:main}, EAAG
consistently outperforms all fixed-direction baselines across
8 environments (Table~\ref{tab:winloss}). The two exceptions occur
in ceiling-effect environments where EAAG already approaches the
oracle. EAAG Pareto-dominates the strongest calibration-based
baseline in most shared environments while requiring no calibration
data. Statistical significance details are in
Appendix~\ref{app:additional}.

Importantly, EAAG's behavior is \emph{qualitatively different}
from baselines: it does not simply improve numbers but exhibits
environment-adaptive gating strategies, as we analyze below.

{\color{blue} \textbf{[DATA NEEDED: Table 3 (tab:main).]} Main results table.
Full method $\times$ environment grid. Rows = methods (base\_only,
always\_trigger, oracle, CaTS, SEAG, CoRefine, CATTS, AUQ, s1\_budget,
BSW, EAAG). Columns = 8 environments, each with SR(\%) and Cost (ro/ep).
Bold best SR per environment. $\dagger$ marks Phase-1-dependent methods.
Similar to FLARE Table 1 (benchmark $\times$ method $\times$ backbone grid).}

{\color{red} \textbf{[FIGURE NEEDED: Figure 5.]} Pareto frontier plot.
x-axis = Cost (rollouts/episode), y-axis = SR(\%).
One panel per environment (or 4 representative ones).
Plot each method as a point. EAAG should be on or near the Pareto frontier.
Similar to FLARE Figure 3 (performance-budget trade-off curves).}

\paragraph{Cross-environment patterns.}
The pattern of results aligns with the Two-Source Model's predictions
along three axes. First, the gap between EAAG and fixed-direction
baselines is largest where direction mismatch is most severe: in
Type~I environments, all ``high uncertainty $\to$ trigger'' baselines
degrade below the no-trigger baseline because the assumed direction
is wrong, and no amount of calibration rescues them. Second, EAAG's
advantage is smallest where direction matters least: in
high-headroom environments where even imprecise gating works, most
methods approach the oracle. Third, EAAG adapts not only direction
but \emph{magnitude} of gating to the environment's rollout value
structure. In rollout-safe environments (TWExpress,
$\Delta{=}+$31.8\,pp), EAAG learns aggressive triggering
(RR=73\%). In rollout-harmful environments (Plancraft,
$\Delta{=}-$7.0\,pp), trigger rate decays over steps
(49\% at step~0 $\to$ $<$20\% later), becoming increasingly
conservative as evidence of rollout harm accumulates. This full
range, from ``always trigger'' to ``almost never trigger,'' emerges
from direction learning alone without explicit headroom estimation.


\subsection{Ablation Studies}
\label{sec:ablation}

The main results show \emph{that} EAAG works; we now ask
\emph{why}. We ablate three components in decreasing order of
expected importance: direction learning, multi-signal features, and
LLM reasoning.

\textbf{Direction is the primary determinant of gate quality.}
Intentionally reversing the learned direction (BSW ablation) causes
catastrophic failure across environments, with an MLP gate using the
wrong direction falling \emph{below} the no-trigger baseline. This
confirms the central prediction of Proposition~\ref{prop:necessity}:
direction is not a minor calibration detail but the primary
determinant of gate quality. The magnitude of degradation correlates
with signal strength (Appendix~\ref{app:proofs}), as the Two-Source
Model predicts.

\textbf{LLM reasoning provides robustness, not marginal accuracy.}
Removing the LLM reasoning step changes SR by ${<}$1\,pp in the
majority of environments. Two environments show larger effects:
one where LLM-generated features capture step-0 criticality, and one
where LLM features introduce noise. The LLM's value is not marginal
accuracy improvement but \emph{generalizability}: it enables
zero-shot feature generation for unseen environments where
task-specific signals cannot be anticipated by universal features
alone.

{\color{red} \textbf{[FIGURE NEEDED: Figure 6.]} Environment-adaptive trigger rates.
Bar chart: x-axis = 8 environments (ordered by rollout headroom $\Delta$),
y-axis = EAAG trigger rate (\%). Show how trigger rate naturally correlates
with headroom: high in TWExpress/HotpotQA, moderate in APPS, low in WebShop,
near-zero in Plancraft. No explicit headroom estimation needed.
Similar to FLARE Figure 2 which shows dynamics across environments.}

\textbf{Gating magnitude emerges from direction learning alone.}
The learned gate adapts trigger rate across environments without
explicit headroom estimation: aggressive triggering in rollout-safe
environments, moderate in mixed environments, and step-dependent
decay in rollout-harmful ones. These environment-adaptive gating
patterns arise from simple logistic regression; the LASSO
coefficients implicitly encode environment-specific triggering
patterns through the learned signal--utility relationship.


\subsection{Two-Source Model Verification}
\label{sec:theory-verification}

The ablations above confirm EAAG's practical advantage; we now test
whether the \emph{theoretical mechanism}, the Two-Source Model
(\S\ref{sec:toy-model}), correctly predicts the observed patterns.
Following \citet{schaeffer2023emergent}, we derive three testable
predictions from the model and verify each against held-out data:
\begin{enumerate}[nosep,label=\textbf{P\arabic*:}]
\item \textbf{Temporal dynamics.} Within the same environment,
  $\rho$ should \emph{decrease} over the episode: early steps mix
  Type~I and Type~D uncertainty, while late steps, after
  information-gathering attempts, isolate the residual Type~I
  component.
\item \textbf{Cross-environment consistency.} Environments with
  similar task structure should exhibit similar $\rho$
  (e.g., FEVER $\approx$ HotpotQA, both search-based QA).
\item \textbf{Signal identity alignment.} In Type~I environments,
  the strongest signal should measure information sufficiency; in
  Type~D environments, decision complexity.
\end{enumerate}

\paragraph{P1: Temporal dynamics (confirmed).}
We split trajectory data into early (step $\leq$ median) and late
(step $>$ median) subsets. In HotpotQA, $\rho$ shifts from $-0.176$
(early) to $-0.437$ (late): late-step entropy reflects deeper
information poverty after failed evidence retrieval. In APPS Intro,
$\rho$ shifts from $+0.102$ (early, non-significant) to $-0.144$
(late, $p{=}0.04$): early exploration diversity gives way to
debugging frustration. In WebShop, the strong positive early-phase
signal ($\rho{=}{+}0.285$) vanishes in the late phase. The pattern
is consistent across all environments with sufficient signal (full
data in Appendix~\ref{app:model}): early steps mix both uncertainty
sources, while late steps isolate the dominant type.

\paragraph{P2: Cross-environment consistency (confirmed).}
We compare same-family pairs. FEVER ($\rho{=}{-}0.119$) and HotpotQA
($\rho{=}{-}0.041$), both search-based QA tasks requiring evidence
retrieval, share negative entropy--utility direction.
APPS Intro ($\rho{=}{+}0.012$) and APPS Interview
($\rho{=}{+}0.317$), both code generation tasks, share non-negative
direction. The magnitude difference within each pair reflects the
degree of Type~I contamination; FEVER has more early-step
information poverty than HotpotQA.

\paragraph{P3: Signal identity alignment (confirmed).}
We examine the LASSO-selected features per environment.
In Type~I environments (HotpotQA, FEVER), the
strongest signal (\texttt{step\_count}, $|\rho|{=}0.494$ and $0.619$)
measures trajectory progress, a proxy for information accumulation.
In WebShop (mixed, choice-dominated),
\texttt{num\_available\_actions} ($\rho{=}{+}0.444$) measures the
decision space size, a proxy for decision complexity. All three
predictions are confirmed, providing converging evidence for the
Two-Source Model.

As a supplementary test, we manipulate information availability
within HotpotQA itself (InfoPoor vs.\ InfoRich conditions) and
confirm that the identity of the most informative signal shifts
accordingly, supporting Observation~2. Full results are reported in
Appendix~\ref{app:controlled}.


\subsection{Robustness of Direction Reversal}
\label{sec:robustness}

A natural concern is whether the observed direction reversal is a
genuine phenomenon or an artifact of confounders. We address this
through three complementary analyses, plus multi-backbone
verification (Appendix~\ref{app:multi-backbone}) showing that
direction depends on the \emph{environment $\times$ model}
interaction, further ruling out model-specific artifacts.

\paragraph{Stratified analysis: controlling for trajectory length.}
Trajectory length varies across environments and correlates with both
entropy and utility. To rule out length as a confounder, we stratify
by step count and recompute $\rho(\text{entropy}, U)$ within each
stratum.
\emph{Result}: Direction reversal persists within trajectory-length
strata where utility variance is non-zero. In HotpotQA, $\rho$ remains
negative across all three strata ($-0.18$/$-0.46$/$-0.42$).
The reversal is not an artifact of trajectory length.

\paragraph{Interventional evidence.}
The BSW ablation (\S\ref{sec:ablation}) provides
\emph{interventional} evidence: reversing the gate's learned
direction causes catastrophic SR degradation across environments,
with the MLP gate falling below the no-trigger baseline. If
direction reversal were merely correlational, this manipulation
should have no systematic effect. The observed damage confirms
that direction encodes a causal relationship.

\paragraph{Gate complexity ablation: direction $\gg$ capacity.}
To verify that the bottleneck is direction rather than gate
complexity, we compare gate architectures with correct vs.\ wrong
direction (Table~\ref{tab:capacity}).

\begin{table}[t]
\caption{Gate complexity ablation on HotpotQA. Direction matters
more than model capacity: a logistic gate with correct direction
outperforms an MLP with wrong direction by 49.9\,pp.}
\label{tab:capacity}
\centering\small
\begin{tabular}{lccc}
\toprule
Gate & Direction & SR (\%) & $\Delta$ \\
\midrule
Logistic (5 feat) & Correct & 95.2 & +46.2 \\
MLP (5 feat)      & Correct & $\sim$95 & $\sim$+46 \\
Hidden state (2560-d) & Correct & $\sim$95 & $\sim$+46 \\
\midrule
Logistic (5 feat) & Wrong & 62.0 & +13.0 \\
MLP (5 feat)      & Wrong & 45.3 & $-$3.7 \\
\midrule
No gate (base)    & --- & 49.0 & 0 \\
\bottomrule
\end{tabular}
\end{table}

With correct direction, all architectures achieve $\sim$95\% SR;
the jump from logistic to MLP to hidden-state probe adds $<$1\,pp.
With wrong direction, higher capacity makes performance \emph{worse}
(MLP 45.3\% $<$ logistic 62.0\%), because the more powerful gate more
precisely targets harmful states. The information hierarchy is:
\textbf{signal direction $\gg$ signal count $\gg$ gate complexity}.
